% Mục 6 theo L0: Tranh luận, kết luận, hướng phát triển (spec §13).
\section{Tranh luận, kết luận và hướng phát triển}
\phutrach{A, B (6.1--6.3) · C (6.4--6.5)}

\subsection{Luật DJ so với điểm mix học (A, B)}
Cùng một cách đo trên MPD và last.fm (\#57, notebook \texttt{11\_lastfm\_mix}; cặp liền nhau ở val, điểm mix học huấn luyện lại trên cặp train của từng bộ):
\begin{table}[htbp]\centering
\caption{Người nghe nối bài so với ngẫu nhiên, và điểm mix học so với điểm mix theo luật (\#57).}
\begin{tabular}{lrr}
\toprule
& MPD & last.fm \\
\midrule
Cặp mix được: thật / ngẫu nhiên (\%) & 4,78 / 3,36 & 4,64 / 3,20 \\
Chênh (điểm \%), CI 95\% theo phiên & +1,42 [1,34; 1,50] & +1,44 [1,42; 1,46] \\
Recall@10 trong batch: học / luật / ngẫu nhiên & 0,037 / 0,019 / 0,010 & 0,096 / 0,021 / 0,010 \\
Spearman giữa điểm học và điểm luật & 0,165 & 0,200 \\
\bottomrule
\end{tabular}
\end{table}
\canviet{A, B: diễn giải; vì sao hai điểm tương quan thấp}

\subsection{Giá của độ mượt (B)}
\canviet{B: NDCG@10 mất bao nhiêu để tỉ lệ cặp mix được tăng bao nhiêu, tại $\alpha$ đã chọn}

\subsection{Vì sao NDCG không đo hết giá trị của chuyển bài mượt (B)}
\canviet{B: chỉ một bài đích; nhiều bài hợp gu và mix mượt bị tính là sai}

\subsection{Kết luận (C)}
\canviet{C}

\subsection{Giới hạn và hướng phát triển (C)}
\canviet{C: đường cong năng lượng cả phiên, cá nhân hóa theo người dùng dài hạn, phát âm thanh thật, khảo sát người nghe}
